\hnheader[Leistungsfähig heißt nicht vertrauenswürdig: fünf Kritikpunkte an ML.][FGSM: $x_{\mathrm{adv}}=x+\varepsilon\,\mathrm{sign}(\nabla_xJ)$]{Kritische}{Perspektiven}{Adversarial Examples, Interpretierbarkeit, Kosten \& Ethik}
\hnsrc{materials/critiques-ml-aut19.pdf (C.\ Chute, T.\ Ding, A.\ Kurenkov, 2019) und materials/critiques-ml.pdf (2018)}

\begin{hngrid}
%
\begin{hnp}[raster multicolumn=2,acc=hnorange]{1}{Fünf Themen}
\begin{enumerate}
\item Adversarial Examples
\item Interpretierbarkeit
\item Kosten: Daten \& Rechenleistung
\item Schwächen der Community
\item Ethische Fragen (Fairness, Bias)
\end{enumerate}
\end{hnp}
%
\begin{hnp}[raster multicolumn=2,acc=hnpurple]{2}{Adversarial Examples}
Winzige, für Menschen kaum sichtbare Störungen erzwingen falsche Vorhersagen ($+\varepsilon\cdot$Rauschen: Panda $\to$ Gibbon). Sie \emph{übertragen} sich auf andere Architekturen und Datensätze und lassen sich physisch bauen (z.\,B.\ Verkehrsschilder).
\end{hnp}
%
\begin{hnp}[raster multicolumn=2,acc=hnteal]{3}{Skizze: Bild $+$ Rauschen}
\centering
\begin{tikzpicture}[every node/.style={draw=hnnavy,rounded corners=1.5pt,font=\scriptsize,inner sep=3pt}]
  \node[fill=hnmint] (a) at (0,0) {Bild $x$};
  \node[draw=none,fill=none] at (1.1,0) {$+\,\varepsilon\,\mathrm{sign}(\nabla_xJ)$};
  \node[fill=hnpink] (b) at (3.2,0) {$x_{\mathrm{adv}}$};
  \node[draw=none,font=\tiny] at (0,-.6) {„Panda“};
  \node[draw=none,font=\tiny] at (3.2,-.6) {„Gibbon“};
\end{tikzpicture}
\end{hnp}
%
\begin{hnp}[raster multicolumn=3,acc=hnnavy]{4}{FGSM \& Verteidigungen}
Fast Gradient Sign Method: gehe in Richtung \emph{steigenden} Verlusts.
\begin{pcode}[FGSM]
\kw{input}: $x,y$, Modell $\theta$, Schrittgröße $\varepsilon$ \ \cm{\# $J$: Verlust, $x_{\mathrm{adv}}$: gestörtes Bild}\\
$g\leftarrow\nabla_xJ(\theta,x,y)$\\
$x_{\mathrm{adv}}\leftarrow x+\varepsilon\cdot\mathrm{sign}(g)$
\end{pcode}
\textbf{Abwehr}: adversariales Training (saubere + gestörte Beispiele), Distillation mit Temperatur-Softmax. Weiterhin ein \emph{Wettrüsten} / offenes Problem.
\end{hnp}
%
\begin{hnp}[raster multicolumn=3,acc=hngreen]{5}{Interpretierbarkeit (Lipton)}
Zwei Bedeutungen: \textbf{algorithmische Transparenz} (Modell selbst verständlich: lineare Modelle, kleine Bäume) vs.\ \textbf{Post-hoc-Erklärung} (Saliency Maps, Grad-CAM, Attention, t-SNE). \emph{Fallacy}: ``lineare Modelle interpretierbar, neuronale Netze Blackbox'' verwechselt beides.\\
Fünf Fragen: Vertrauen, Kausalität, Übertragbarkeit, Informativität, Fairness -- Metriken wie Accuracy/AUROC erfassen sie nicht (Vertiefung: Kapitel SHAP und Fairness).
\end{hnp}
%
\begin{hnex}[raster multicolumn=3]{Beispiel: warum lineare Modelle anfällig sind}
$f(x)=w\T x$, $x\in\R^{1000}$, $|w_i|=0.01$. Störung $\eta=\varepsilon\,\mathrm{sign}(w)$ mit $\varepsilon=0.1$ pro Feature.\\
\hnsol\ $w\T\tilde x=w\T x+w\T\eta$ mit $w\T\eta=\varepsilon\|w\|_1=0.1\cdot1000\cdot0.01=\ora{1.0}$.\\
Jedes Feature ändert sich nur um $0.1$, die Ausgabe aber um $1.0$ $\Rightarrow$ in \emph{hohen Dimensionen} summieren sich kleine Störungen (Goodfellow).
\end{hnex}
%
\begin{hnp}[raster multicolumn=3,acc=hnrose]{6}{Kosten \& Gesellschaft}
\begin{itemize}
\item \textbf{Daten \& Compute}: moderne Modelle brauchen riesige Datenmengen und Rechenleistung -- Zugang und Umwelt.
\item \textbf{Bias in Daten} $\Rightarrow$ ungerechte Entscheidungen gegenüber unterrepräsentierten Gruppen.
\item \textbf{Verantwortung}: Evaluation über reine Testfehler hinaus.
\end{itemize}
\end{hnp}
%
\begin{hnp}[raster multicolumn=2,acc=hngreen]{7}{Was tun?}
\begin{itemize}
\item Robustheit testen
\item Erklärungen bereitstellen
\item Fairness-Metriken prüfen
\end{itemize}
\end{hnp}
%
\begin{hnp}[raster multicolumn=2,acc=hnred]{8}{Offene Probleme}
\begin{hncon}
\item robuste Verteidigung
\item Kausalität statt Korrelation
\item Übertragung auf neue Verteilungen
\end{hncon}
\end{hnp}
%
\begin{hnp}[raster multicolumn=2,acc=hnorange]{9}{Key Takeaways}
\begin{hnchk}
\item FGSM: $x+\varepsilon\,\mathrm{sign}(\nabla_xJ)$
\item Zwei Arten von Interpretierbarkeit
\item Genauigkeit $\ne$ Vertrauen
\end{hnchk}
\end{hnp}
%
\begin{hnp}[raster multicolumn=6,acc=hnpurple,colback=hnlilac!30]{?}{Selbsttest: kannst du das erklären?}
\hnquiz{Was ist FGSM?}{$x_{\mathrm{adv}}=x+\varepsilon\,\mathrm{sign}(\nabla_xJ)$: Schritt in Richtung steigenden Verlusts.}{Zwei Arten von Interpretierbarkeit?}{Algorithmische Transparenz vs.\ Post-hoc-Erklärung.}{Warum sind lineare Modelle anfällig?}{$w\T\eta=\varepsilon\|w\|_1$ wächst mit der Dimension.}
\end{hnp}
%
\begin{hnbanner}[raster multicolumn=6]
„Ein Modell ist erst dann gut, wenn man ihm auch dann vertrauen kann, wenn niemand hinsieht.“
\end{hnbanner}
\end{hngrid}
